% The load-store queue: loads and stores in program order between the commit
% pointer and the issue pointer.  The load in LSQ6 searches the earlier
% entries, skips loads and non-matching stores, and takes the value of the
% nearest store with the same address.
\begin{tikzpicture}[x=1cm,y=1cm, font=\footnotesize\sffamily, >={Stealth[length=1.8mm]},
  hd/.style={font=\footnotesize\bfseries\sffamily},
  used/.style={draw, fill=ln-accent!12, minimum height=0.45cm, inner sep=0pt},
  free/.style={draw, fill=white, minimum height=0.45cm, inner sep=0pt},
  me/.style={draw, fill=ln-accent!35, minimum height=0.45cm, inner sep=0pt},
  lab/.style={font=\scriptsize\sffamily, text=ln-muted},
  ann/.style={font=\scriptsize\sffamily, anchor=west}]
  % column headings
  \node[hd] at (0.35,0.45) {L/S};
  \node[hd] at (1.25,0.45) {Addr};
  \node[hd] at (2.25,0.45) {Value};
  \node[hd] at (2.95,0.45) {C};
  % rows: entry / type / address / value / style
  \foreach \r/\t/\a/\v/\st in {1/{}/{}/{}/free, 2/S/0x100/17/used,
      3/L/0x108/4/used, 4/S/0x108/25/used, 5/L/{}/{}/used,
      6/L/0x100/{}/me, 7/S/{}/{}/used, 8/{}/{}/{}/free} {
    \pgfmathsetmacro{\y}{-(\r-1)*0.45}
    \node[\st, minimum width=0.7cm] at (0.35,\y) {\t};
    \node[\st, minimum width=1.1cm, font=\scriptsize\ttfamily] at (1.25,\y) {\a};
    \node[\st, minimum width=0.9cm] at (2.25,\y) {\v};
    \node[\st, minimum width=0.5cm] at (2.95,\y) {};
    \node[lab, anchor=east] at (-0.05,\y) {LSQ\r};
  }
  % completion bits: set for LSQ2-LSQ4, and for LSQ6 by the forwarding
  \foreach \y in {-0.45,-0.9,-1.35} {
    \draw[thick] (2.83,\y) -- (2.92,\y-0.09) -- (3.08,\y+0.1);
  }
  \draw[thick, ln-accent] (2.83,-2.25) -- (2.92,-2.34) -- (3.08,-2.15);
  % forwarded value in the searching load
  \node[text=ln-accent, font=\footnotesize\bfseries\sffamily] at (2.25,-2.25) {17};
  % pointers
  \draw[->, thick, ln-accent] (-1.45,-0.45) -- (-0.85,-0.45);
  \node[font=\scriptsize\sffamily, text=ln-accent, anchor=east] at (-1.45,-0.45)
    {\iflangja{コミット}{commit}};
  \draw[->, thick, ln-accent] (-1.45,-3.15) -- (-0.85,-3.15);
  \node[font=\scriptsize\sffamily, text=ln-accent, anchor=east] at (-1.45,-3.15)
    {\iflangja{発行}{issue}};
  % search of the load in LSQ6
  \draw[thick, ln-accent] (3.25,-2.25) -- (3.6,-2.25) -- (3.6,-0.45);
  \draw[->, thick, ln-accent] (3.6,-0.45) -- (3.25,-0.45);
  \foreach \y in {-0.9,-1.35,-1.8} { \draw[ln-accent] (3.52,\y) -- (3.68,\y); }
  \node[ann, text=ln-accent] at (3.8,-2.25)
    {\iflangja{アドレス 0x100 のロード：前方を探索}{load at 0x100: searches earlier entries}};
  \node[ann, text=ln-muted] at (3.8,-1.8) {\iflangja{ロード：無視}{load: skipped}};
  \node[ann, text=ln-muted] at (3.8,-1.35) {\iflangja{ストア，0x108：不一致}{store, 0x108: no match}};
  \node[ann, text=ln-muted] at (3.8,-0.9) {\iflangja{ロード：無視}{load: skipped}};
  \node[ann, text=ln-accent] at (3.8,-0.45)
    {\iflangja{ストア，0x100：一致，値 17 を転送}{store, 0x100: match, value 17 forwarded}};
  \node[ann, text=ln-muted] at (3.8,-2.7)
    {\iflangja{後続のストア：探索しない}{later store: not examined}};
  % program order
  \draw[->, ln-muted] (-2.55,0.2) -- (-2.55,-3.35);
  \node[lab, rotate=90] at (-2.8,-1.57) {\iflangja{プログラム順}{program order}};
\end{tikzpicture}
